\documentclass[12pt,a4paper]{article}
\usepackage[a4paper,margin=2.25cm,headheight=15pt]{geometry}
\usepackage[T1]{fontenc}
\usepackage{iftex}
\ifPDFTeX
\usepackage[utf8]{inputenc}
\fi
\usepackage{lmodern}
\usepackage{graphicx}
\usepackage{booktabs,longtable,tabularx,array}
\usepackage{enumitem}
\usepackage{xcolor}
\usepackage{fancyhdr}
\usepackage{float}
\usepackage{xurl}
\usepackage{hyperref}
\setlength{\emergencystretch}{2em}
\hypersetup{colorlinks=true,linkcolor=blue,urlcolor=blue}
\pagestyle{fancy}
\fancyhf{}\lhead{Francanglais Compiler}\rhead{SDD}\cfoot{\thepage}
\newcommand{\placeholder}[1]{\textcolor{red}{\textbf{#1}}}
\newcommand{\diagram}[2]{\begin{figure}[H]\centering\includegraphics[width=#2\textwidth,height=0.82\textheight,keepaspectratio]{diagrams/#1.png}\caption{#1}\end{figure}}

\title{\textbf{Software Design Description}\\[0.4em]
\Large Francanglais Compiler and Dataset Platform}
\author{Group members and matricules:\\[0.5em]
\placeholder{Name 1 --- Matricule 1}\\
\placeholder{Name 2 --- Matricule 2}\\
\placeholder{Name 3 --- Matricule 3}}
\date{Version 1.0 --- \today}

\begin{document}
\maketitle
\thispagestyle{empty}
\begin{center}
\textit{Group work document. Replace all red placeholders before submission.}
\end{center}
\newpage
\tableofcontents
\newpage

\section{Design Basis and Evidence}
The design is grounded in the source-controlled implementation:
\texttt{data\_collector/App.py}, \texttt{dataset.py}, and
\texttt{audio\_utils.py}. The collector uses a CustomTkinter tabbed interface,
CSV persistence, and a managed audio directory. The repository also includes
a compiled frontend bundle and local ignored backend/compiler artifacts; those
are represented as integration boundaries, not as verified public APIs.

\section{Architectural Overview}
The system follows a local layered design:
\begin{enumerate}
\item \textbf{Presentation layer}: \texttt{App} builds Collect, Browse \& Edit,
and Stats tabs and handles user commands.
\item \textbf{Dataset layer}: \texttt{dataset.py} owns CSV initialization,
reading, append, replacement, update, deletion, duplicate checks, and counts.
\item \textbf{Audio layer}: \texttt{audio\_utils.py} records, copies, saves,
and plays audio with graceful capability detection.
\item \textbf{Compiler boundary}: a future lexer/classifier/parser consumes
dataset evidence without coupling the UI directly to parsing internals.
\end{enumerate}
\diagram{component}{0.98}

\section{Deployment View}
The implemented collector is deployed on one contributor workstation with
Python runtime, local CSV storage, and local audio files. A compiler service
or API adapter is optional and is shown separately to prevent confusion
between current and planned deployment.
\diagram{deployment}{0.98}

\section{Detailed Module Design}
\subsection{App}
\texttt{App} initializes the dataset file, configures the dark UI, creates
three tabs, and owns transient form state such as the pending audio filename
and selected browse row. \texttt{save\_entry} reads the form, creates an ID
and timestamp, calls \texttt{dataset.append\_entry}, and refreshes visible
state. Browse operations delegate persistence to dataset helpers. Audio
controls delegate recording and playback to \texttt{audio\_utils}.

\subsection{Dataset}
The dataset module defines the canonical field list, entry types, and
categories. \nolinkurl{ensure_dataset_file} establishes the CSV header;
\texttt{load\_all} returns dictionaries; \texttt{append\_entry} appends one
row; \texttt{save\_all} rewrites the complete file for updates/deletes.
\texttt{text\_exists} normalizes case and whitespace for duplicate detection.

\subsection{Audio utilities}
\texttt{AUDIO\_RECORDING\_AVAILABLE} is determined by optional imports.
\texttt{Recorder} buffers microphone frames and returns one NumPy array.
\texttt{save\_recording} writes a UUID-named WAV file. \texttt{attach\_file}
copies an existing file to a UUID-named destination. \texttt{play\_audio}
tries sounddevice/soundfile first and then the operating-system default
player; it returns false when the path is missing or both paths fail.

\subsection{Compiler boundary}
The compiler design is intentionally staged: tokenize text, classify tokens
using exact/fuzzy and learned lexical evidence, parse a mixed-language
sentence with a grammar, and return a structured analysis. This boundary
should be implemented behind an adapter so the collector remains usable
without the compiler and so compiler analysis cannot silently rewrite data.

\subsection{Class relationships and multiplicities}
The class diagram includes both application-defined classes, \texttt{App}
and \nolinkurl{audio_utils.Recorder}, together with the relevant external
classes \nolinkurl{customtkinter.CTk} and \nolinkurl{sounddevice.InputStream}.
It shows selected attributes and methods rather than every GUI widget.
\texttt{App} inherits from \texttt{CTk} and owns the optional recorder;
the audio helper module does not own a recorder instance.

The \texttt{module}, \texttt{record}, and \texttt{logical storage}
stereotypes distinguish supporting design elements from implemented
classes. In particular, \texttt{dataset.py} and \texttt{audio\_utils.py}
are modules whose functions are shown as static operations.
\texttt{DatasetEntry} describes the CSV row/Python dictionary schema;
\texttt{DatasetFile} and \texttt{AudioFile} describe stored files.
These three names are not Python model classes.

\diagram{class-domain}{0.98}

\begin{longtable}{>{\raggedright\arraybackslash}p{4.6cm}>{\raggedright\arraybackslash}p{9.8cm}}
\toprule
\textbf{Relationship} & \textbf{Multiplicity and meaning}\\ \midrule
\endhead
App--Recorder & One app owns zero or one recorder, depending on audio capability. Each recorder created by the collector belongs to one app.\\
Recorder--InputStream & One recorder owns zero or one stream. The stream is created when recording starts and closed and released when it stops.\\
DatasetFile--DatasetEntry & One CSV file contains zero or more rows; each stored row belongs to one CSV file in this data model.\\
DatasetEntry--AudioFile & A row optionally references one audio file by filename. A file may have zero or more references because filenames are not constrained to be unique across rows.\\
\bottomrule
\end{longtable}

Filled diamonds denote composition and sit at the owning (whole) end.
The audio-file link is an association, not composition: deleting a CSV row
does not delete its audio file. This logical link does not guarantee that
the referenced file still exists. Inheritance and function-call dependencies
do not have multiplicities; adding cardinalities to those arrows would not
be valid UML.

\section{Behavioral Design}
\subsection{Use cases}
\diagram{use-case}{0.98}

\subsection{Entry collection activity}
\diagram{activity-collect}{0.90}

\subsection{Save sequence}
The sequence shows an optional existing-file attachment before saving,
followed by mutually exclusive empty-text and valid-text paths. An empty
submission returns a warning without a CSV write. On the valid path,
the application persists the entry, refreshes counts and recent entries,
and clears the form while retaining the contributor.

Every synchronous message starts an activation on its receiving lifeline,
and every dashed reply ends that activation. A self-call creates a nested
activation; the caller remains active while waiting. The contributor's
activation spans each user interaction. Module lifelines represent executing
functions, and the file-I/O lifeline abstracts local filesystem operations,
not an additional application class or remote service. Individual widget
calls and lower-level file helpers are omitted. Live recording and storage
exceptions are outside this particular scenario.
\diagram{sequence-collect}{0.98}

\subsection{Entry lifecycle}
\diagram{state-entry}{0.90}

\section{Interfaces and Contracts}
\begin{longtable}{>{\raggedright\arraybackslash}p{6.8cm}>{\raggedright\arraybackslash}p{7.8cm}}
\toprule
\textbf{Interface} & \textbf{Contract}\\ \midrule
\endhead
\nolinkurl{dataset.append_entry(entry)} & Receives a dictionary containing the canonical fields and appends one UTF-8 CSV row.\\
\nolinkurl{dataset.update_entry(id, fields)} & Loads rows, updates the matching ID, and rewrites the CSV.\\
\nolinkurl{dataset.delete_entry(id)} & Loads rows, removes matching ID, and rewrites the CSV.\\
\nolinkurl{audio_utils.attach_file(source, dir)} & Copies a source media file to the managed directory and returns its generated filename.\\
\nolinkurl{audio_utils.play_audio(path)} & Returns true only when playback was successfully started.\\
\nolinkurl{CompilerAdapter.analyse(text)} & Proposed contract returning tokens, token labels, parse status, and explanatory diagnostics.\\
\bottomrule
\end{longtable}

\section{Persistence and Error Handling}
The CSV header is created before reads or writes. Update and delete use
read-modify-write semantics, which is suitable for a single desktop user but
not for concurrent writers. Empty text is rejected at the UI boundary.
Missing audio paths and playback failures are surfaced as false results and,
for pending playback, a visible error dialog. Optional recording import
failure disables recording while preserving attachment and data-management
features.

\section{Security and Privacy Design}
Local environment files may contain credentials and are excluded from the
documentation and Git history. Audio and linguistic entries can contain
personal or culturally sensitive information; access should be restricted to
the project team, contributor consent should be obtained, and exports should
exclude secrets. A future service adapter must validate input, authenticate
maintainer operations, and avoid logging raw audio or credentials.

\section{Verification Strategy}
\begin{itemize}
\item Unit-test CSV initialization, UTF-8 round trips, append, update, delete,
duplicate normalization, counts, and empty-input validation.
\item Test audio capability-off behavior without a microphone, plus attach and
missing-file playback paths.
\item Test compiler tokenization and parsing against representative
French/English/Pidgin mixtures and malformed input.
\item Perform an end-to-end manual scenario covering Collect, Browse \& Edit,
Stats, and audio playback.
\item Regenerate diagrams from the PlantUML sources so the PNGs remain
consistent with the design text.
\end{itemize}

\section{Design Decisions and Future Work}
\begin{itemize}
\item CSV is retained for transparent, portable Phase 1 collection.
\item The compiler remains decoupled until the corpus and grammar are stable.
\item A database, schema migration strategy, authentication, and service API
are future work for multi-user deployment.
\item The illustrative grammar in the coursework JSON requires linguistic
validation before being promoted to a production parser.
\end{itemize}

\end{document}
